\documentclass{article}%
\usepackage[T1]{fontenc}%
\usepackage[utf8]{inputenc}%
\usepackage{lmodern}%
\usepackage{textcomp}%
\usepackage{lastpage}%
\usepackage{geometry}%
\geometry{margin=1in,headheight=14pt,headsep=25pt}%
\usepackage{hyperref}%
\usepackage{enumitem}%
\usepackage{fancyhdr}%
\usepackage{xcolor}%
\usepackage{url}%
\usepackage{breakurl}%
%

            \hypersetup{
                colorlinks=true,
                linkcolor=blue,
                filecolor=magenta,
                urlcolor=blue
            }
            \pagestyle{fancy}
            \fancyhf{}
            \rhead{Generated on \today}
            \cfoot{\thepage}
            
            % Configure enumeration settings
            \setlist[enumerate,1]{label=\arabic*.}
            \setlist[enumerate,2]{label=\alph*.}
            \setlist[enumerate,3]{label=\Alph*.}
            \setlist[enumerate]{itemsep=0.5em}
            
            % Define a command for red text
            \newcommand{\incorrect}[1]{\textcolor{red}{#1}}
        %
\lhead{2024{-}10{-}28{-}Farkas}%
\title{Summary for 2024{-}10{-}28{-}Farkas}%
\author{Generated by Scribe.AI}%
\date{\today}%
%
\begin{document}%
\normalsize%
\maketitle%
\section{Summary}%
\label{sec:Summary}%
Farkas' Lemma is a fundamental theorem in linear programming that establishes relationships between the feasibility of a system of linear inequalities and the existence of solutions to a related dual system.  The lecture explored various formulations of this lemma, highlighting its versatility in analyzing linear programs.

*   **Farkas' Lemma:** A system of linear inequalities, Ax ≤ b, has no solution if and only if there exists a non-negative vector y such that Aᵀy = 0 and bᵀy < 0.

*   **Alternative Formulations:**  The lemma can be expressed using equality constraints (Ax = b) and non-negativity constraints on x, leading to alternative conditions for solvability involving a dual system in y.

*   **Three Variants:**  The lecture presented three equivalent variants of Farkas' Lemma, each addressing the solvability of systems with different combinations of equality and inequality constraints, and non-negativity conditions on the solution vector x.

*   **Connection to Fredholm Alternatives:** The lecture highlighted the relationship between Farkas' Lemma and the Fredholm Alternatives, demonstrating its broader mathematical significance beyond linear programming.

*   **Conditions for Non-Existence of Solutions:**  Several slides provided conditions under which the primal systems (Ax ≤ b or Ax = b) have no solution, expressed in terms of the existence of a solution to a corresponding dual system.

*   **Geometric Interpretation:** The proof of Farkas' Lemma was outlined using a geometric approach, involving the projection of a point onto a cone defined by the matrix A.

*   **Related Theorems:** The lecture included a problem requiring the derivation of several theorems (Gordan, Farkas, Stiemke, Ville, and Tucker) closely related to Farkas' Lemma, showcasing its wide-ranging applications in mathematical analysis.

%
\end{document}